%LTeX: language=fr-FR
\documentclass{article}
\usepackage{tasks}

\usepackage{amsmath,amssymb,amsfonts}
\usepackage{graphicx}
\usepackage{amsthm}
\usepackage[french]{babel}
\usepackage{enumitem}
\usepackage{currfile}

\usepackage{tikz}
\usetikzlibrary{calc,intersections}
\usepackage{qrcode}
\usepackage{mathrsfs}
\usepackage{enumitem}
\usepackage[french]{babel}
\usepackage{tcolorbox}
\usepackage{fancyhdr}
\usepackage{multicol}
\usepackage[headsep=3mm,footskip=5mm,includefoot,includehead]{geometry}
%showframe
\usepackage[lastpage,user]{zref}

\newtheorem*{theorem*}{Théorème}

\setlength{\parindent}{0cm}

\newenvironment{clist}{
  \begin{center}
  \begin{inparaenum}}{
    \end{inparaenum}
\end{center}}

\pagestyle{empty}
\graphicspath{ {../../Images/} }
\input{\currfiledir ../def.tex}
\def \Chapitre {Chapitre 3 -- Généralités sur les fonctions}
\def \ChapitreCourt {Ch. 3 -- Généralités sur les fonctions}


\geometry{left=1.5cm, right=1.5cm, top=1cm, bottom=1cm}
\setlength{\headheight}{14.0pt}

\theoremstyle{definition}
\newtheorem{exercice}{Exercice}

\begin{document}
\pagestyle{fancy}

\fancyhf{} % clear existing header/footer entries
% We don't need to specify the O coordinate
\fancyhead[L]{\large \textbf{\Niveau}}
\fancyhead[C]{\large \textbf{\Chapitre}}

\large

Au cycle 4, les élèves ont découvert la notion de fonction et ses différents modes de représentation : expression algébrique, tableau de valeurs, représentation graphique, programme de calcul. Ils connaissent le vocabulaire de base : variable, fonction, image, antécédent et la notation $f(x)$. En seconde, on consolide la notion de fonction, comme exprimant la dépendance d'une variable par rapport à une autre, en passant d'un registre à l'autre (algébrique et graphique).

\begin{paragraph}{Pré-requis}
  \begin{itemize}
    \item Cycle 4 : notion de fonction, image, antécédent, notation $f(x)$.
    \item Calcul d'une expression littérale pour une valeur donnée (priorités, nombres relatifs).
    \item Coordonnées d'un point dans un repère.
    \item Chapitre 2 : intervalles, réunion d'intervalles.
  \end{itemize}
\end{paragraph}

\begin{paragraph}{Contenu}
    \begin{itemize}
        \item Fonction à valeurs réelles définie sur un intervalle ou une réunion finie d'intervalles de $\mathbb{R}$. Ensemble de définition.
        \item Image, antécédent. Fonction donnée par un tableau de valeurs, une formule, une courbe.
        \item Courbe représentative : la courbe d'équation $y = f(x)$ est l'ensemble des points du plan dont les coordonnées $(x\,;\,y)$ vérifient $y = f(x)$.
    \end{itemize}
\end{paragraph}

\begin{paragraph}{Objectifs}
  \begin{itemize}
    \item Calculer une image à partir d'une formule ; lire une image et des antécédents sur un tableau ou une courbe.
    \item Exploiter l'équation $y = f(x)$ d'une courbe : appartenance d'un point, calcul de coordonnées.
    \item Résoudre graphiquement une équation ou une inéquation du type $f(x) = k$, $f(x) < k$.
    \item Résoudre graphiquement une équation ou une inéquation du type $f(x) = g(x)$, $f(x) < g(x)$.
    \item Modéliser par des fonctions des situations issues des mathématiques (problèmes d'aires) ou de la vie courante.
  \end{itemize}
\end{paragraph}

\begin{paragraph}{Notions reportées}
    \begin{itemize}
        \item Signe d'une fonction, tableau de signes, variations et extremums : chapitre \og Fonctions : variations et signes \fg{} (décembre).
        \item Fonctions de référence, parité : au fil de l'année.
    \end{itemize}
\end{paragraph}

\begin{paragraph}{Histoire des mathématiques}
    \begin{itemize}
        \item La lente élaboration de la notion de fonction : Newton, Leibniz, Euler, jusqu'à la définition de Dirichlet ; importance de la notation $f(x)$.
    \end{itemize}
\end{paragraph}

\end{document}
